\section{Introduction}\label{sec:intro}

\subsection{Exactly divergence-free elements}
Most mixed finite element methods for incompressible flow impose the constraint $\div u=0$ only weakly, against the discrete pressures. The velocity error then depends on the pressure, and this can be large when pressure gradients are large or the viscosity is small; the review of John, Linke, Merdon, Neilan and Rebholz \cite{JohnEtAl2017} explains why this matters and what can be done about it. One remedy is to choose the velocity and pressure spaces so that the divergence of every discrete velocity is itself a discrete pressure. Then the weak constraint forces $\div u_h=0$ pointwise, mass is conserved exactly, and on fitted meshes with strongly imposed no-slip data the velocity error does not see the pressure.

The oldest such pair is due to Scott and Vogelius \cite{ScottVogelius1985}: continuous piecewise-$P_k$ velocities with discontinuous piecewise-$P_{k-1}$ pressures. In two dimensions they proved $\div\Vh=\Pih$ for $k\ge4$ on meshes without singular vertices, and Guzm\'an and Scott \cite{GuzmanScott2019} later made the stability constant explicit in the local mesh geometry. Lower degrees become possible on split meshes; for example, quadratic velocities and linear pressures are stable on Clough--Tocher splits \cite{ArnoldQin1992}.

Three dimensions are much harder, and the theory is still incomplete. Zhang proved stability of the Scott--Vogelius pair on Alfeld (barycentric) refinements for $k\ge3$ \cite{Zhang2005}, and on Freudenthal meshes for $k\ge6$ \cite{Zhang2011}; Guzm\'an and Neilan \cite{GuzmanNeilan2018} gave a short proof on barycentric refinements in every dimension $d$, for $k\ge d$, built from a \emph{local} stability result on a single split simplex. Fu, Guzm\'an and Neilan \cite{FuGuzmanNeilan2020} constructed exact smooth piecewise-polynomial sequences on Alfeld splits, with commuting projections. On Worsey--Farin splits, Guzm\'an, Lischke and Neilan \cite{GuzmanLischkeNeilan2022} and Fabien, Guzm\'an, Neilan and Zytoon \cite{FabienGuzmanNeilanZytoon2022} obtained divergence-free pairs of low degree, whose pressure spaces carry extra conditions at the singular edges of the split. Neilan \cite{Neilan2015} constructs smooth discrete de Rham complexes in three dimensions, and Farrell, Mitchell and Scott \cite{FarrellMitchellScott2024} formulate conjectures on the Stokes complex on Freudenthal meshes; we know of no unconditional stability result for the Scott--Vogelius pair on general unsplit tetrahedral meshes. Neilan's review \cite{Neilan2020} surveys the field. In this paper, all unconditional results are on Alfeld refinements.

\subsection{Curved walls}
Exact incompressibility interacts badly with curved boundaries. Replacing a smooth domain by an inscribed polygon or polyhedron is the cheapest way to mesh it. For scalar second-order problems this costs order $h^{3/2}$ in $H^1$ \cite{BergerScottStrang1972,Scott1975}, and classical remedies either correct the boundary values \cite{BrambleDupontThomee1972,DupontGuzmanScott2025} or fit the boundary with curved elements. For Scott--Vogelius elements the cost is higher. In two dimensions, a smooth divergence-free field that vanishes on two edges meeting at a polygon vertex has zero gradient there, and discrete divergence-free fields inherit such \emph{zero-gradient} constraints at the vertices of an inscribed wall. If no-slip is imposed strongly on the inscribed polygon, the $H^1$ error can then lock at order $h^{1/2}$ \cite{ScottPrize,GjerdeScott2024}; in the companion paper \cite{PadhiLong2D} we trace the lock to wall vertices that lie in only two triangles.

Three routes avoid the problem. (i) \emph{Fit the geometry.} Neilan and Otus \cite{NeilanOtus2021} proved optimal convergence of Scott--Vogelius elements on isoparametric curved triangulations in two dimensions, Durst and Neilan \cite{DurstNeilan2024} treated general degree on smooth planar domains, and Chen and Liu \cite{ChenLiu2026} have very recently extended the isoparametric approach to three dimensions. (ii) \emph{Correct the boundary condition} on a mesh that does not reach the true boundary: Liu, Neilan and Otus \cite{LiuNeilanOtus2023} combine a divergence-free method with a Taylor boundary correction in two dimensions. Boundary-value corrections go back to Bramble, Dupont and Thom\'ee \cite{BrambleDupontThomee1972}; related ideas underlie the shifted boundary method \cite{MainScovazzi2018} and cut methods with boundary-value correction \cite{BurmanHansboLarson2018,BurmanHansboLarson2019Stokes}. (iii) \emph{Impose the wall condition weakly.} Nitsche's method \cite{Nitsche1971,Stenberg1995,FreundStenberg1995,JuntunenStenberg2009} and its penalty and penalty-free relatives \cite{Babuska1973,Burman2012,BoiveauBurman2016} impose Dirichlet data through a consistent boundary term, and are the standard tool of unfitted and cut finite element methods for Stokes flow \cite{BurmanHansbo2014}, including recent exactly divergence-free cut methods \cite{LiuNeilanOlshanskii2023,FrachonNilssonZahedi2024,BurmanHansboLarson2024}. Gjerde and Scott have used Nitsche's method for slip conditions in Navier--Stokes flow \cite{GjerdeScott2022}.

\subsection{Scott's zero-gradient prize and the two-dimensional answer}
Gjerde and Scott \cite{GjerdeScott2024} computed drag on a cylinder with Scott--Vogelius elements. They kept the inscribed polygon but imposed no-slip by Nitsche's method, with a penalty $\mu/h$ and no pressure term in the boundary form, and observed errors far below those of strong imposition. Scott's zero-gradient prize \cite{ScottPrize} asks for a proof that this Scott--Vogelius--Nitsche method converges, with the error $\hG^{3/2}+h^k$ one expects from the scalar theory ($\hG$ is the size of the wall faces and $h$ the bulk mesh size), and adds: \emph{if this does not hold, why not?}

The two-dimensional answer is given in the companion paper \cite{PadhiLong2D}, and its central mechanism in a short letter \cite{CavalcantePadhi2026}. The mechanism is simple. Nitsche's boundary form comes from integrating the full stress $\sigma(u,p)n=D(u)n-pn$ by parts. The method as printed in \cite[(27)--(28)]{GjerdeScott2024} keeps a viscous flux term, $\partial_nu$, with its symmetric counterpart and the penalty, but has no pressure term $-pn$. At the wall, nothing in the discrete momentum balance then opposes the pressure except the penalty $(\mu/h)\,u_h\cdot n$, so the discrete fluid leaks through the wall with normal velocity approximately $(h/\mu)(p-\pbar)$. The wall mean $\pbar$ is subtracted because an exactly divergence-free velocity has zero net flux through the wall. The printed method therefore converges only like $(h/\mu)^{1/2}$ in energy and $h/\mu$ in $H^1$ unless the wall pressure is constant. Adding back the plain traction $\ip{p_h}{v\cdot n}$ makes the discrete system singular, because constant pressures no longer see the divergence; adding the traction with the wall mean of the discrete pressure removed gives a well-posed, exactly divergence-free method, which we call $\CNS$. Its velocity coincides with that of the zero-net-flux formulation of Frachon, Nilsson and Zahedi \cite{FrachonNilssonZahedi2024}; the device is theirs. In two dimensions $\CNS$ attains $\hG^{3/2}+h^k$, the term $\hG^{3/2}$ is sharp, and the Nitsche penalty must grow with the ratio of bulk to boundary mesh size \cite{PadhiLong2D}. The constant-pressure benchmark of \cite{GjerdeScott2024,ScottPrize}, for which the leak vanishes, was analysed independently by Cavalcante \cite{Cavalcante}.

\subsection{This paper}
We answer the same question in three dimensions. The wall $\Gamma$ is a smooth closed surface, possibly with several components and not necessarily convex (the exterior of the unit ball in a box is the model case). It is replaced by a closed triangulated surface $\Gh$ whose vertices lie on $\Gamma$, and the region between $\Gh$ and a polyhedral outer boundary is meshed by straight-sided tetrahedra. The velocity is continuous piecewise $P_k$, the pressure discontinuous piecewise $P_{k-1}$, and the wall condition is imposed by Nitsche's method. We study the printed method, which we call $\GS$, and the corrected method $\CNS$.

\paragraph{Where $\hG^{3/2}$ comes from.} A wall face $F$ with vertices on $\Gamma$ is a chord of the surface: it lies at a small distance $d$ from $\Gamma$, and $d$ vanishes at the vertices. To leading order $d$ is the quadratic bump
\[
d\approx\tfrac12\sum_{i<j}\hh(e_{ij},e_{ij})\,\mu_i\mu_j
\]
in the barycentric coordinates $\mu_i$ of $F$, where $e_{ij}$ are the edge vectors and $\hh$ is the second fundamental form of $\Gamma$ (Lemma~\ref{geom:face}). The exact velocity vanishes on $\Gamma$, so on $\Gh$ it is approximately $d\,\partial_Nu$, a field of size $\hG^2$ that oscillates on the scale $\hG$. Such a trace has $H^{1/2}$-norm of order $\hG^{3/2}$, and this is the geometric term in the error. The lower bound asks the converse question: can the discrete divergence-free fields near a face see the bump? The bump is described by three numbers per face, the values of $\hh$ on the three edges, so the lower bound reduces to three \emph{edge moments} of the normal derivative on each face. The weight that results is $|\hh|\,\partial_Nu$: on flat parts of the wall the faces lie in $\Gamma$, there is no bump, and no $\hG^{3/2}$ term.

\paragraph{What is new in three dimensions.} Four things change, and they account for most of the length of the paper.
\begin{itemize}
\item \emph{Curvature enters through the second fundamental form.} In two dimensions the sagitta of a chord is a single number times the curvature. On a triangle it is a quadratic form, and the lower bound needs discrete fields that respond to all three edge weights. On Alfeld refinements with $k\ge4$ one macro-element suffices; for $k=3$ it provably does not, and the witnesses must be built on the star of an interior edge.
\item \emph{A face-mean defect.} In two dimensions the tilt between the chord normal and the normal of $\Gamma$ is odd about the chord midpoint, which is also the centre of the chord's circumcircle, and this gives a useful cancellation. On a triangle inscribed in a sphere the face mean of the tilt is proportional to the offset between centroid and circumcentre (Lemma~\ref{geom:sag}), which is nonzero on every non-equilateral face. We show that this defect changes constants but no rate.
\item \emph{No stream functions.} The two-dimensional constructions use stream functions. We replace them by a Piola--Hestenes extension of the flux $2$-form of the velocity (Lemma~\ref{geom:ext}) and by interpolate-and-correct constructions based on the discrete inf-sup condition.
\item \emph{Geometry and stability.} In surface finite elements the closest-point projection $\Gh\to\Gamma$ is assumed to be bijective; we derive it from the volume mesh. Stability of the Scott--Vogelius pair in three dimensions is known only on split meshes, and it must hold uniformly over the family of perturbed domains. On Worsey--Farin splits the discrete inf-sup condition fails for the full pressure space, as is known \cite{GuzmanLischkeNeilan2022,FabienGuzmanNeilanZytoon2022}; this is the three-dimensional counterpart of the singular vertices of two dimensions.
\end{itemize}

\subsection{Main results}
We state the results informally; precise statements, with all hypotheses, are in the sections cited. Throughout, $h$ is the bulk mesh size that appears in Scott's penalty $\mu/h$, $\hG$ is the largest diameter of a wall face, $\gamma:=\mu\hG/h$ is the effective wall penalty, $\rho$ is the ratio of $h$ to the smallest wall-face diameter, $\pbar$ is the mean of $p$ over $\Gamma$ and $G:=\norm{p-\pbar}_{L^2(\Gamma)}$ measures how far the wall pressure is from constant. The exact solution is assumed smooth up to the boundary; at the edges and corners of the outer box this is a hypothesis on the solution, not a consequence of smooth data (Remark~\ref{set:rem:corner}).

\begin{enumerate}[label=\textbf{\arabic*.},leftmargin=1.6em]
\item \emph{Geometry.} If the tetrahedral mesh is shape regular and the discrete domain does not reach the far side of the obstacle (condition (OB) of Section~\ref{sec:setting}), then the closest-point projection maps $\Gh$ homeomorphically onto $\Gamma$ and $\Gh$ is a normal graph over $\Gamma$ of height $O(\hG^2)$ (Theorems~\ref{geom:orient} and~\ref{geom:mesh}). Neither hypothesis can be dropped (Examples~\ref{ex:sliver} and~\ref{ex:filled}).
\item \emph{Stability.} On Alfeld refinements with $k\ge3$ the discrete inf-sup constant is bounded below uniformly in $h$ (Lemma~\ref{st:IS3}). On meshes whose wall faces are subdivided, such as Worsey--Farin splits, the inf-sup condition fails for the full pressure space (Proposition~\ref{st:wf}), as is known.
\item \emph{The printed method leaks.} If the wall pressure is not constant, $\GS$ converges with error exactly of order $(h/\mu)^{1/2}G$ in energy and $h/\mu$ in $H^1$. To leading order the discrete velocity satisfies the leak condition $(\mu/h)\,u_h\cdot n\approx p-\pbar$ on $\Gh$, with constant exactly $1$ in the slip and energy norms (Theorems~\ref{er:A}, \ref{er:C}, \ref{er:B}, Corollary~\ref{er:Bp}).
\item \emph{The leak limit.} For every penalty above the coercivity threshold, the scaled error $(\mu/h)(\ut-u_h)$, minus its geometric part, converges in $H^1$ to the Stokes flow $U$ driven by the Dirichlet data $-(p-\pbar)n$ on $\Gamma$, at rate $\hG^{1/2}$ (Theorem~\ref{lk:up}). On Alfeld refinements with $k\ge4$ the rate is exactly $\frac12$, with a lower bound weighted by the second fundamental form (Theorem~\ref{lk:low}). So the $H^1$ constant of $\GS$ is $\norm{\nabla U}/G$ (Corollary~\ref{lk:gen}).
\item \emph{The corrected method attains Scott's rate.} $\CNS$ is well posed and exactly divergence-free, and for bounded $\gamma\ge\gamma_2$ its $H^1$ error is at most $C(\hG^{3/2}+h^k)$ provided $\hG\ge h^2$ (Theorem~\ref{er:D}); the proviso disappears for outer data with zero discrete net flux. On Alfeld refinements this holds for every $k\ge3$ with no further hypothesis.
\item \emph{Sharpness.} For $\CNS$, and for $\GS$ with any wall pressure,
\[
\norm{\nabla(\ut-u_h)}\ \ge\ c\,\hG^{3/2}\,\bigl\|\,|\hh|\,\partial_Nu\bigr\|_{L^2(\Gamma)}-C\hG^2,
\]
on Alfeld refinements for every $k\ge3$ and on every shape-regular mesh for $k\ge9$ (Theorems~\ref{sh:lb}, \ref{sh:alfeld}, \ref{sh:k3main}, Proposition~\ref{sh:bubble}). So the $\hG^{3/2}$ term of Scott's rate cannot be improved whenever the wall shear is nonzero somewhere on a curved part of the wall.
\item \emph{Drag.} The drag, lift and torque of $\GS$ have the error $-(h/\mu)K_\psi$ to leading order, with $K_\psi$ given by a reciprocity formula (Theorem~\ref{dr:gs}); those of $\CNS$ converge at rate $\frac32$, and at rate $2$ modulo a regularity assumption for a dual problem at the outer boundary (Theorem~\ref{dr:cns}).
\item \emph{Penalty.} A penalty $\mu\ge C\rho$ suffices for coercivity. Conversely, coercivity on the divergence-free space fails for $\mu$ below a constant times $\rho$, under explicit conditions at a smallest wall face: either its Alfeld macro-element has a certified positive single-cell constant (exact certificates for a range of shapes), or the mesh is locally quasi-uniform there (a continuum witness) (Theorems~\ref{pen:single}, \ref{pen:cont}).
\item \emph{Navier--Stokes.} Under discrete linearised-stability conditions, which hold for small data and, for small $h$, on nonsingular solution branches, the upper bounds, the leak law, the leak limit and the drag statements carry over to steady Navier--Stokes flow (Corollary~\ref{ns:cor}, Theorem~\ref{ns:oseen}).
\item \emph{Computations.} A first three-dimensional implementation confirms exact incompressibility, the optimal rates on a manufactured problem and the leak law with constant tending to $1$; the sphere tests we could run are still pre-asymptotic for the rate $\hG^{3/2}$ (Section~\ref{sec:numerics}).
\end{enumerate}
The constants of the lower bounds are explicit on Alfeld refinements with $k\ge4$ (Appendix~\ref{app:constants}). Table~\ref{tab:summary} collects the results with their scope and their status.

\begin{table}[t]
\centering\footnotesize
\caption{Summary of results. ``Exact'' marks proofs that use a finite computation in exact rational or interval arithmetic, reproducible from Appendix~\ref{app:repro}. (IS3) is the uniform discrete inf-sup condition; it is a theorem on Alfeld refinements with $k\ge3$. The other named hypotheses are explained in Table~\ref{tab:hyp}.}\label{tab:summary}
\begin{tabular}{>{\raggedright\arraybackslash}p{0.27\textwidth}>{\raggedright\arraybackslash}p{0.27\textwidth}>{\raggedright\arraybackslash}p{0.23\textwidth}>{\raggedright\arraybackslash}p{0.13\textwidth}}
\toprule
Statement & Size & Scope & Status\\
\midrule
$\pi:\Gh\to\Gamma$ homeomorphism & normal graph, height $O(\hG^2)$ & shape regular $+$ (OB) & proved\\
uniform discrete inf-sup & $\beta$ independent of $h$ & Alfeld, $k\ge3$ & proved\\
$\GS$, energy error & $\asymp(h/\mu)^{1/2}G$ & (IS3) & proved\\
$\GS$, $H^1$ error & $\asymp h/\mu$, constant $\norm{\nabla U}/G$ & (IS3), all $\gamma\ge\gamma_0$ & proved\\
leak limit & rate exactly $\hG^{1/2}$ & upper: (IS3); lower: Alfeld, $k\ge4$ & proved; lower bound exact\\
$\CNS$, $H^1$ error & $\le C(\hG^{3/2}+h^k)$ & (IS3), bounded $\gamma\ge\gamma_2$ & proved\\
lower bound $c\,\hG^{3/2}$ & weight $|\hh|\,\partial_Nu$ & Alfeld $k\ge3$; any mesh $k\ge9$ & proved; exact\\
$\GS$ drag, lift, torque & error $-(h/\mu)K_\psi+\dots$ & (IS3) & proved\\
$\CNS$ drag, lift, torque & rate $\frac32$; rate $2$ & (IS3) & proved; rate $2$ modulo (E$_Z$)\\
penalty threshold $\mu^\ast$ & $c\rho\le\mu^\ast\le C\rho$ & lower: certified cell or (LQ) & proved; single cell exact\\
Navier--Stokes & as for Stokes & (L), (L$^\ast$), (L$_D$), (L$^\dagger$) & conditional\\
Worsey--Farin & inf-sup fails for full $\Pih$ & subdivided wall faces & known; exact count\\
computations & rates on a box, leak law & $k=3,4$ & numerical\\
\bottomrule
\end{tabular}
\end{table}

\subsection{Related work and novelty}\label{sec:novelty}
We list what is known, what we take from others, and what we believe is new.

\paragraph{Known, and used.} Stability of the Scott--Vogelius pair on Alfeld refinements is due to Zhang \cite{Zhang2005} and Guzm\'an and Neilan \cite{GuzmanNeilan2018}; our uniform-in-$h$ version (Lemma~\ref{st:IS3}) combines the local result of \cite[Theorem~3.1]{GuzmanNeilan2018}, Bernardi--Raugel face fluxes \cite{BernardiRaugel1985} and the continuity of the continuous inf-sup constant under Lipschitz perturbations of the domain due to Bernardi, Costabel, Dauge and Girault \cite{BCDG2016}. Chen and Liu use the same combination for polyhedral approximations of curved domains \cite[Lemma~3.9]{ChenLiu2026}, and our argument follows theirs in structure; we claim no novelty for it. The mean-free traction of $\CNS$ is the zero-net-flux device of Frachon, Nilsson and Zahedi \cite{FrachonNilssonZahedi2024}. The failure of the inf-sup condition for the full pressure space on Worsey--Farin splits is known \cite{GuzmanLischkeNeilan2022,FabienGuzmanNeilanZytoon2022}; we only record a short proof and an exact count for our setting. The order $h^{3/2}$ for inscribed polygonal boundaries is classical for scalar problems \cite{BergerScottStrang1972,Scott1975}. The structure of our lift theorem (a proper local homeomorphism with one sheet) is that of Amenta, Choi, Dey and Leekha \cite{AmentaChoiDeyLeekha2000}.

\paragraph{Closest recent work.} Chen and Liu \cite{ChenLiu2026} prove the optimal, pressure-robust rate $h^k$ for Piola-mapped Scott--Vogelius elements with $k\ge3$ on Alfeld splits of \emph{isoparametric} curved tetrahedral meshes, with strongly imposed Dirichlet data and a commuting-interpolant load. That is the right method if one is willing to build curved elements. Our setting is the cheap one that Scott's question is about: straight-sided, affinely mapped tetrahedra, an inscribed triangulated wall, standard Nitsche terms, and no curved element maps, multipliers or boundary corrections. Its best possible rate is $\hG^{3/2}+h^k$, and our contribution is to show that $\CNS$ attains this rate and that it cannot be improved; we do not compete with \cite{ChenLiu2026} on rate. Liu, Neilan and Otus \cite{LiuNeilanOtus2023} reach order $h^k$ in two dimensions on a mesh inside the curved domain with a mean-free boundary multiplier, a zero-flux constraint and a Taylor boundary correction; the divergence-free cut methods \cite{LiuNeilanOlshanskii2023,FrachonNilssonZahedi2024,BurmanHansboLarson2024} are unfitted, and their analyses concern that setting. None of these works treats the Nitsche method on inscribed polyhedra, the leak of the printed method, or lower bounds. Cavalcante \cite{Cavalcante} analyses the two-dimensional constant-pressure benchmark. The two-dimensional versions of our results are in \cite{PadhiLong2D}.

\paragraph{New, to our knowledge.}
\begin{enumerate}[label=(\alph*)]
\item The lift theorem: the closest-point projection $\Gh\to\Gamma$ is a homeomorphism as a \emph{consequence} of shape regularity of the volume mesh and the global condition (OB), with examples showing that neither can be dropped (Section~\ref{sec:geometry}).
\item The three-dimensional leak law of the printed method, with exact constants, and the leak limit uniformly in the penalty, with the exact rate $\hG^{1/2}$ and a lower bound weighted by the second fundamental form (Sections~\ref{sec:error} and~\ref{sec:leak}).
\item The analysis of $\CNS$ on inscribed polyhedra: well-posedness and the rate $\hG^{3/2}+h^k$, unconditional on Alfeld refinements for $k\ge3$.
\item The sharp lower bound $c\,\hG^{3/2}\norm{|\hh|\,\partial_Nu}$, with the reduction to three edge moments per face, the exact witnesses for $k\ge4$, the proof that one macro-element cannot work for $k=3$, and the edge-star construction for $k=3$, including saddle and parabolic points (Section~\ref{sec:sharp} and Appendix~\ref{app:k3}).
\item Drag, lift and torque errors of both methods in three dimensions, and the Navier--Stokes and Oseen versions.
\item Necessity of a penalty proportional to $\rho$ in three dimensions, with exact single-cell certificates and a continuum witness; explicit constants for the lower bounds.
\end{enumerate}
Items (b)--(f) have two-dimensional counterparts in \cite{PadhiLong2D}. Where a two-dimensional proof carries over with only notational changes, we give the three-dimensional replacement of each dimension-dependent ingredient and refer to \cite{PadhiLong2D} for the unchanged algebra.

\subsection{Organisation and conventions}
Section~\ref{sec:setting} fixes the setting and collects the named hypotheses in Table~\ref{tab:hyp}. Section~\ref{sec:geometry} studies inscribed surfaces and proves the lift theorem, Section~\ref{sec:stability} the stability results, Section~\ref{sec:error} the error analysis of both methods, Section~\ref{sec:sharp} the lower bounds, Section~\ref{sec:leak} the leak limit and drag, Section~\ref{sec:penalty} the penalty threshold, Section~\ref{sec:ns} the Navier--Stokes extension and Section~\ref{sec:numerics} the computations. Section~\ref{sec:open} discusses what remains open. Appendix~\ref{app:k3} contains the degree-three witnesses, Appendix~\ref{app:constants} the explicit constants and the details of the computer-assisted certificates, and Appendix~\ref{app:repro} the command behind every computed number.

A statement marked \status{computer-assisted, exact} relies on a finite computation in exact rational arithmetic (or, where stated, interval arithmetic) whose output is a certificate that can be rechecked independently; \status{numerical} marks floating-point evidence that is not part of any proof. Constants $C,c$ never depend on $h,\hG,\mu,\gamma,\rho$ unless this is displayed.
